\subsection{Orthogonal Projections}

Throughout this subsection, $\mathbb{F}$ is $\mathbb{R}$ or $\mathbb{C}$ and $\mathcal{V}$ is an inner-product space over $\mathbb{F}$.

\begin{proposition} \label{prop:orthogonal-projection-onto-line} \hlt{Orthogonal Projection onto a Line}\\
Let $\mathbf{u},\mathbf{v}\in\mathcal{V}$ with $\mathbf{v}\neq\mathbf{0}$, and set $c\defeq\frac{\innerprod{\mathbf{u},\mathbf{v}}}{\norm{\mathbf{v}}^2}$ and $\mathbf{w}\defeq\mathbf{u}-c\mathbf{v}$. Then $\mathbf{u}=c\mathbf{v}+\mathbf{w}$ and $\mathbf{w}\perp\mathbf{v}$.\\
The vector $\operatorname{proj}_{\mathbf{v}}(\mathbf{u})\defeq c\mathbf{v}=\innerprod{\mathbf{u},\frac{\mathbf{v}}{\norm{\mathbf{v}}}}\frac{\mathbf{v}}{\norm{\mathbf{v}}}$ is the \hlt{orthogonal projection} of $\mathbf{u}$ onto $\operatorname{span}(\mathbf{v})$.
\begin{proof}
Since $\mathbf{v}\neq\mathbf{0}$, $\norm{\mathbf{v}}^2>0$ by definiteness, so $c$ is well-defined. The identity $\mathbf{w}=\mathbf{u}-c\mathbf{v}$ implies $\mathbf{u}=c\mathbf{v}+\mathbf{w}$.\\
By linearity, $\innerprod{\mathbf{w},\mathbf{v}}=\innerprod{\mathbf{u},\mathbf{v}}-c\innerprod{\mathbf{v},\mathbf{v}}=\innerprod{\mathbf{u},\mathbf{v}}-\frac{\innerprod{\mathbf{u},\mathbf{v}}}{\norm{\mathbf{v}}^2}\norm{\mathbf{v}}^2=0$, so $\mathbf{w}\perp\mathbf{v}$.\\
Since $\norm{\mathbf{v}}\in\mathbb{R}_{>0}$, $\innerprod{\mathbf{u},\frac{\mathbf{v}}{\norm{\mathbf{v}}}}\frac{\mathbf{v}}{\norm{\mathbf{v}}}=\frac{\innerprod{\mathbf{u},\mathbf{v}}}{\norm{\mathbf{v}}^2}\mathbf{v}=c\mathbf{v}$ by Proposition \ref{prop:basic-inner-product-identities}(iv).\\
\hlt{Structural Commentary}: The chosen coefficient makes the residual orthogonal.\\
\hlt{Key Insight}: Projection isolates the parallel component and leaves an orthogonal residual.
\end{proof}
\end{proposition}

Figure \ref{fig:orthogonal-decomposition-one-vector} depicts the decomposition $\mathbf{u}=c\mathbf{v}+\mathbf{w}$ with $c\mathbf{v}\in\operatorname{span}(\mathbf{v})$ and $\mathbf{w}\perp\mathbf{v}$.

\begin{figure}[H]
  \centering
  \begin{tikzpicture}[
  x=0.88cm,
  y=0.88cm,
  >=Stealth,
  every node/.style={font=\small},
  axis/.style={gray!55, line width=0.45pt},
  vec/.style={->, blue!75!black, line width=0.9pt},
  comp/.style={->, teal!70!black, line width=0.85pt},
  orth/.style={->, red!70!black, line width=0.85pt},
  point/.style={circle, fill=black, inner sep=1.1pt}
]
  \coordinate (O) at (0,0);
  \coordinate (P) at (2.8,0);
  \coordinate (U) at (2.8,1.35);
  \coordinate (V) at (1.25,0);

  \draw[axis] (-0.25,0) -- (3.65,0) node[right] {$\operatorname{span}(\mathbf{v})$};
  \draw[vec] (O) -- (U) node[midway, above left] {$\mathbf{u}$};
  \draw[comp] (O) -- (P) node[midway, below] {$c\mathbf{v}$};
  \draw[orth] (P) -- (U) node[midway, right] {$\mathbf{w}$};
  \draw[vec, blue!45!black] (O) -- (V) node[midway, below] {$\mathbf{v}$};

  \draw[gray!60] (2.58,0) -- (2.58,0.22) -- (2.8,0.22);
  \node[point, label=below left:{$\mathbf{0}$}] at (O) {};
  \node[point] at (P) {};
  \node[point] at (U) {};
  \node[above] at (3.0,1.35) {$\mathbf{u}=c\mathbf{v}+\mathbf{w}$};
\end{tikzpicture}

  \caption{Writing $\mathbf{u}$ as a component parallel to $\mathbf{v}$ plus a component orthogonal to $\mathbf{v}$.}
  \label{fig:orthogonal-decomposition-one-vector}
\end{figure}

\begin{definition} \label{def:orthogonal-projection-subspace} \hlt{Orthogonal Projection onto a Subspace}\\
Let $\mathcal{U}$ be finite-dimensional subspace of $\mathcal{V}$. By Theorem \ref{thm:orthogonal-direct-sum-decomposition}, every $\mathbf{v}\in\mathcal{V}$ has unique decomposition $\mathbf{v}=\mathbf{u}+\mathbf{w}$ with $\mathbf{u}\in\mathcal{U}$, $\mathbf{w}\in\mathcal{U}^{\perp}$.\\
The \hlt{orthogonal projection} of $\mathcal{V}$ onto $\mathcal{U}$ is the map $P_{\mathcal{U}}:\mathcal{V}\to\mathcal{V}$ defined by $P_{\mathcal{U}}\mathbf{v}\defeq\mathbf{u}$.
\end{definition}

\par Figure \ref{fig:orthogonal-projection-subspace} illustrates the $\mathcal{U}$-component $P_{\mathcal{U}}\mathbf{v}$ and the residual $\mathbf{v}-P_{\mathcal{U}}\mathbf{v}\in\mathcal{U}^{\perp}$.

\begin{figure}[H]
  \centering
  \begin{tikzpicture}[
  x=0.88cm,
  y=0.88cm,
  >=Stealth,
  every node/.style={font=\small},
  subspace/.style={blue!65!black, line width=0.8pt},
  vec/.style={->, blue!75!black, line width=0.95pt},
  proj/.style={->, teal!70!black, line width=0.9pt},
  residual/.style={->, red!70!black, line width=0.85pt},
  point/.style={circle, fill=black, inner sep=1pt}
]
  \coordinate (O) at (0,0);
  \coordinate (A) at (-0.45,-0.11);
  \coordinate (B) at (4.6,1.15);
  \coordinate (P) at (2.4,0.6);
  \coordinate (V) at (1.95,2.4);

  \draw[subspace] (A) -- (B) node[right] {$\mathcal{U}$};
  \draw[vec] (O) -- (V) node[pos=0.46, left=7pt] {$\mathbf{v}$};
  \draw[proj] (O) -- (P) node[pos=0.58, below=6pt] {$P_{\mathcal{U}}\mathbf{v}$};
  \draw[residual] (P) -- (V) node[pos=0.5, right=5pt] {$\mathbf{v}-P_{\mathcal{U}}\mathbf{v}$};
  \draw[gray!60] (2.356,0.775) -- (2.526,0.817) -- (2.57,0.642);

  \node[point, label=below left:{$\mathbf{0}$}] at (O) {};
  \node[point] at (P) {};
  \node[point, label=above:{$\mathbf{v}$}] at (V) {};
\end{tikzpicture}

  \caption{Orthogonal projection decomposes $\mathbf{v}$ into $P_{\mathcal{U}}\mathbf{v}$ plus an orthogonal residual.}
  \label{fig:orthogonal-projection-subspace}
\end{figure}

\begin{proposition} \label{prop:orthogonal-projection-properties} \hlt{Properties of Orthogonal Projection}\\
Let $\mathcal{U}$ be a finite-dimensional subspace of $\mathcal{V}$. Then:
\begin{enumerate}[label=(\roman*)]
\setlength{\itemsep}{0pt}
\item $P_{\mathcal{U}}\in\LL(\mathcal{V})$;
\item $(\forall\,\mathbf{u}\in\mathcal{U})\,P_{\mathcal{U}}\mathbf{u}=\mathbf{u}$;
\item $(\forall\,\mathbf{w}\in\mathcal{U}^{\perp})\,P_{\mathcal{U}}\mathbf{w}=\mathbf{0}$;
\item $\operatorname{range}(P_{\mathcal{U}})=\mathcal{U}$;
\item $(\forall\,\mathbf{v}\in\mathcal{V})\,\mathbf{v}-P_{\mathcal{U}}\mathbf{v}\in\mathcal{U}^{\perp}$;
\item $\ker(P_{\mathcal{U}})=\mathcal{U}^{\perp}$;
\item $P_{\mathcal{U}}^2=P_{\mathcal{U}}$;
\item $(\forall\,\mathbf{v}\in\mathcal{V})\,\norm{P_{\mathcal{U}}\mathbf{v}}\leq\norm{\mathbf{v}}$;
\item if $(\mathbf{e}_1,\ldots,\mathbf{e}_m)$ is an orthonormal basis of $\mathcal{U}$, then $(\forall\,\mathbf{v}\in\mathcal{V})\,P_{\mathcal{U}}\mathbf{v}=\sum\limits_{j=1}^{m}\innerprod{\mathbf{v},\mathbf{e}_j}\mathbf{e}_j$.
\end{enumerate}
\begin{proof}
\begin{enumerate}[label=(\roman*)]
\setlength{\itemsep}{0pt}
\item Fix $a\in\mathbb{F}$ and $\mathbf{v}_1,\mathbf{v}_2\in\mathcal{V}$. By Theorem \ref{thm:orthogonal-direct-sum-decomposition}, write $\mathbf{v}_r=\mathbf{u}_r+\mathbf{w}_r$ with $\mathbf{u}_r\in\mathcal{U}$ and $\mathbf{w}_r\in\mathcal{U}^{\perp}$ for $r\in\{1,2\}$.\\
Since $\mathcal{U}$ and $\mathcal{U}^{\perp}$ are subspaces, $a\mathbf{v}_1+\mathbf{v}_2=(a\mathbf{u}_1+\mathbf{u}_2)+(a\mathbf{w}_1+\mathbf{w}_2)$ is its orthogonal decomposition.\\
Uniqueness therefore implies $P_{\mathcal{U}}(a\mathbf{v}_1+\mathbf{v}_2)=aP_{\mathcal{U}}\mathbf{v}_1+P_{\mathcal{U}}\mathbf{v}_2$, so $P_{\mathcal{U}}\in\LL(\mathcal{V})$.
\item Fix $\mathbf{u}\in\mathcal{U}$. The decomposition $\mathbf{u}=\mathbf{u}+\mathbf{0}$ has $\mathbf{0}\in\mathcal{U}^{\perp}$ by Proposition \ref{prop:basic-properties-orthogonal-complement}(i), so $P_{\mathcal{U}}\mathbf{u}=\mathbf{u}$.
\item Fix $\mathbf{w}\in\mathcal{U}^{\perp}$. The decomposition $\mathbf{w}=\mathbf{0}+\mathbf{w}$ has $\mathbf{0}\in\mathcal{U}$, so $P_{\mathcal{U}}\mathbf{w}=\mathbf{0}$.
\item The defining decomposition places every value of $P_{\mathcal{U}}$ in $\mathcal{U}$, so $\operatorname{range}(P_{\mathcal{U}})\subseteq\mathcal{U}$. The reverse inclusion holds by (ii).
\item Fix $\mathbf{v}\in\mathcal{V}$. The defining decomposition is $\mathbf{v}=P_{\mathcal{U}}\mathbf{v}+(\mathbf{v}-P_{\mathcal{U}}\mathbf{v})$ with $\mathbf{v}-P_{\mathcal{U}}\mathbf{v}\in\mathcal{U}^{\perp}$.
\item The inclusion $\mathcal{U}^{\perp}\subseteq\ker(P_{\mathcal{U}})$ holds by (iii). Conversely, if $P_{\mathcal{U}}\mathbf{v}=\mathbf{0}$, then $\mathbf{v}=\mathbf{v}-P_{\mathcal{U}}\mathbf{v}\in\mathcal{U}^{\perp}$ by (v).
\item Fix $\mathbf{v}\in\mathcal{V}$. By (iv), $P_{\mathcal{U}}\mathbf{v}\in\mathcal{U}$, so $P_{\mathcal{U}}(P_{\mathcal{U}}\mathbf{v})=P_{\mathcal{U}}\mathbf{v}$ by (ii). Hence $P_{\mathcal{U}}^2=P_{\mathcal{U}}$.
\item Fix $\mathbf{v}\in\mathcal{V}$ and set $\mathbf{p}\defeq P_{\mathcal{U}}\mathbf{v}$ and $\mathbf{r}\defeq\mathbf{v}-\mathbf{p}$. Then $\mathbf{p}\in\mathcal{U}$ by (iv) and $\mathbf{r}\in\mathcal{U}^{\perp}$ by (v).\\
Hence $\norm{\mathbf{v}}^2=\norm{\mathbf{p}}^2+\norm{\mathbf{r}}^2\geq\norm{\mathbf{p}}^2$ by Theorem \ref{thm:pythagorean-inner-product}. Nonnegativity of norms implies $\norm{P_{\mathcal{U}}\mathbf{v}}\leq\norm{\mathbf{v}}$.
\item Fix $\mathbf{v}\in\mathcal{V}$ and set $\mathbf{p}\defeq\sum\limits_{j=1}^{m}\innerprod{\mathbf{v},\mathbf{e}_j}\mathbf{e}_j\in\mathcal{U}$. For $k\in\{1,\ldots,m\}$, $\innerprod{\mathbf{v}-\mathbf{p},\mathbf{e}_k}=\innerprod{\mathbf{v},\mathbf{e}_k}-\sum\limits_{j=1}^{m}\innerprod{\mathbf{v},\mathbf{e}_j}\innerprod{\mathbf{e}_j,\mathbf{e}_k}=0$ by orthonormality. \\
If $\mathbf{x}\in\mathcal{U}$, write $\mathbf{x}=\sum\limits_{k=1}^{m}a_k\mathbf{e}_k$; then $\innerprod{\mathbf{v}-\mathbf{p},\mathbf{x}}=\sum\limits_{k=1}^{m}\overline{a_k}\innerprod{\mathbf{v}-\mathbf{p},\mathbf{e}_k}=0$ by Proposition \ref{prop:basic-inner-product-identities}(iii)--(iv).\\
Thus $\mathbf{v}-\mathbf{p}\in\mathcal{U}^{\perp}$, and $P_{\mathcal{U}}\mathbf{v}=\mathbf{p}$ by uniqueness in Theorem \ref{thm:orthogonal-direct-sum-decomposition}.
\end{enumerate}
\hlt{Structural Commentary}: The orthogonal decomposition organizes all nine arguments.\\
\hlt{Key Insight}: Projection preserves the $\mathcal{U}$-component and annihilates the $\mathcal{U}^{\perp}$-component.
\end{proof}
\end{proposition}

\begin{definition} \label{def:rank-one-projection-matrix} \hlt{Rank-One Projection Matrix}\\
Let $\mathbf{u}=(u_1,\ldots,u_n)^t\in\mathbb{F}^n$ be a unit vector. The \hlt{rank-one projection matrix} onto $\operatorname{span}(\mathbf{u})$ is $P_{\mathbf{u}}\defeq[u_i\overline{u_j}]_{i,j=1}^{n}\in\operatorname{Mat}(n,n,\mathbb{F})$.\\
It satisfies $(\forall\,\mathbf{x}\in\mathbb{F}^n)\,P_{\mathbf{u}}\mathbf{x}=\innerprod[\mathbb{F}^n]{\mathbf{x},\mathbf{u}}\mathbf{u}$.
\end{definition}

\begin{proposition} \label{prop:rank-one-projection-matrix-properties} \hlt{Properties of Rank-One Projection Matrices}\\
Let $\mathbf{u}\in\mathbb{F}^n$ be a unit vector. Then:
\begin{enumerate}[label=(\roman*)]
\setlength{\itemsep}{0pt}
\item $\operatorname{range}(P_{\mathbf{u}})=\operatorname{span}(\mathbf{u})$;
\item $\ker(P_{\mathbf{u}})=\operatorname{span}(\mathbf{u})^{\perp}$;
\item $(\forall,i,j\in\{1,\ldots,n\})\,(P_{\mathbf{u}})_{ij}=\overline{(P_{\mathbf{u}})_{ji}}$;
\item $P_{\mathbf{u}}^2=P_{\mathbf{u}}$.
\end{enumerate}
\begin{proof}
\begin{enumerate}[label=(\roman*)]
\setlength{\itemsep}{0pt}
\item For $\mathbf{x}\in\mathbb{F}^n$, $P_{\mathbf{u}}\mathbf{x}=\innerprod[\mathbb{F}^n]{\mathbf{x},\mathbf{u}}\mathbf{u}\in\operatorname{span}(\mathbf{u})$, so $\operatorname{range}(P_{\mathbf{u}})\subseteq\operatorname{span}(\mathbf{u})$. Since $P_{\mathbf{u}}\mathbf{u}=\innerprod[\mathbb{F}^n]{\mathbf{u},\mathbf{u}}\mathbf{u}=\mathbf{u}$, the reverse inclusion also holds.
\item Fix $\mathbf{x}\in\mathbb{F}^n$. Since $\mathbf{u}\neq\mathbf{0}$, $P_{\mathbf{u}}\mathbf{x}=\innerprod[\mathbb{F}^n]{\mathbf{x},\mathbf{u}}\mathbf{u}=\mathbf{0}$ if and only if $\innerprod[\mathbb{F}^n]{\mathbf{x},\mathbf{u}}=0$, which is equivalent to $\mathbf{x}\in\operatorname{span}(\mathbf{u})^{\perp}$.
\item For $i,j\in\{1,\ldots,n\}$, $(P_{\mathbf{u}})_{ij}=u_i\overline{u_j}=\overline{u_j\overline{u_i}}=\overline{(P_{\mathbf{u}})_{ji}}$.
\item For $i,k\in\{1,\ldots,n\}$, $(P_{\mathbf{u}}^2)_{ik}=\sum\limits_{j=1}^{n}u_i\overline{u_j}u_j\overline{u_k}=u_i\overline{u_k}\sum\limits_{j=1}^{n}\abs{u_j}^2=u_i\overline{u_k}=(P_{\mathbf{u}})_{ik}$ because $\mathbf{u}$ is a unit vector.
\end{enumerate}
\hlt{Structural Commentary}: The projection formula proves all four matrix identities.\\
\hlt{Key Insight}: The matrix preserves the $\mathbf{u}$ direction and annihilates its orthogonal complement.
\end{proof}
\end{proposition}

\begin{theorem} \label{thm:orthogonal-projection-minimizes-distance} \hlt{Orthogonal Projection Minimizes Distance}\\
Let $\mathcal{U}$ be a finite-dimensional subspace of $\mathcal{V}$. If $\mathbf{v}\in\mathcal{V}$ and $\mathbf{u}\in\mathcal{U}$, then $\norm{\mathbf{v}-P_{\mathcal{U}}\mathbf{v}}\leq\norm{\mathbf{v}-\mathbf{u}}$, with equality if and only if $\mathbf{u}=P_{\mathcal{U}}\mathbf{v}$.
\begin{proof}
Set $\mathbf{p}\defeq P_{\mathcal{U}}\mathbf{v}$, $\mathbf{r}\defeq\mathbf{v}-\mathbf{p}$, and $\mathbf{h}\defeq\mathbf{p}-\mathbf{u}$. By Proposition \ref{prop:orthogonal-projection-properties}(iv)--(v), $\mathbf{p},\mathbf{h}\in\mathcal{U}$ and $\mathbf{r}\in\mathcal{U}^{\perp}$, so $\mathbf{r}\perp\mathbf{h}$.\\
Since $\mathbf{v}-\mathbf{u}=\mathbf{r}+\mathbf{h}$, $\norm{\mathbf{v}-\mathbf{u}}^2=\norm{\mathbf{r}}^2+\norm{\mathbf{h}}^2\geq\norm{\mathbf{r}}^2=\norm{\mathbf{v}-P_{\mathcal{U}}\mathbf{v}}^2$ by Theorem \ref{thm:pythagorean-inner-product}. Nonnegativity of norms proves the inequality.
\begin{enumerate}[label=\roman*.]
\setlength{\itemsep}{0pt}
\item[$(\Rightarrow)$] If equality holds, then the preceding identity forces $\norm{\mathbf{h}}=0$, so $\mathbf{h}=\mathbf{0}$ by definiteness and $\mathbf{u}=\mathbf{p}=P_{\mathcal{U}}\mathbf{v}$. \item[$(\Leftarrow)$] If $\mathbf{u}=P_{\mathcal{U}}\mathbf{v}$, then $\mathbf{h}=\mathbf{0}$ and equality holds in the same identity.
\end{enumerate}
\hlt{Structural Commentary}: Pythagoras isolates the nonnegative correction term.\\
\hlt{Key Insight}: The orthogonal projection is the unique nearest point in the subspace.
\end{proof}
\end{theorem}

\begin{definition} \label{def:distance-to-subspace} \hlt{Distance to a Subspace}\\
Let $\mathcal{U}$ be a nonempty subset of $\mathcal{V}$. The \hlt{distance} from $\mathbf{v}\in\mathcal{V}$ to $\mathcal{U}$ is $d(\mathbf{v},\mathcal{U})\defeq\inf_{\mathbf{u}\in\mathcal{U}}\norm{\mathbf{v}-\mathbf{u}}$.\\
If $\mathcal{U}$ is a finite-dimensional subspace, then $d(\mathbf{v},\mathcal{U})=\norm{\mathbf{v}-P_{\mathcal{U}}\mathbf{v}}$ by Theorem \ref{thm:orthogonal-projection-minimizes-distance}.
\end{definition}

\begin{definition} \label{def:fourier-coefficients-and-expansion} \hlt{Fourier Coefficients and Expansion}\\
Let $\mathcal{O}=(\mathbf{e}_1,\ldots,\mathbf{e}_m)$ be orthonormal list, $\mathbf{v}\in\mathcal{V}$. The \hlt{Fourier coefficients} of $\mathbf{v}$ relative to $\mathcal{O}$ are $c_j(\mathbf{v})\defeq\innerprod{\mathbf{v},\mathbf{e}_j}$ for $j\in\{1,\ldots,m\}$.\\
The \hlt{Fourier expansion} of $\mathbf{v}$ onto $\mathcal{O}$ is $\widehat{\mathbf{v}}_{\mathcal{O}}\defeq\sum\limits_{j=1}^{m}c_j(\mathbf{v})\mathbf{e}_j$.
\end{definition}

\begin{theorem} \label{thm:fourier-expansion-best-approximation-finite} \hlt{Fourier Expansion onto a Finite Orthonormal Set}\\
Let $\mathcal{O}=(\mathbf{e}_1,\ldots,\mathbf{e}_m)$ be orthonormal, let $\mathcal{S}=\operatorname{span}(\mathcal{O})$, and let $\mathbf{v}\in\mathcal{V}$. Then:
\begin{enumerate}[label=(\roman*)]
\setlength{\itemsep}{0pt}
\item $\widehat{\mathbf{v}}_{\mathcal{O}}$ is the unique $\mathbf{s}\in\mathcal{S}$ with $\mathbf{v}-\mathbf{s}\in\mathcal{S}^{\perp}$;
\item $(\forall\,\mathbf{s}\in\mathcal{S})\,\norm{\mathbf{v}-\widehat{\mathbf{v}}_{\mathcal{O}}}\leq\norm{\mathbf{v}-\mathbf{s}}$, and $\widehat{\mathbf{v}}_{\mathcal{O}}$ is the only $\mathbf{s}\in\mathcal{S}$ for which equality holds;
\item $\norm{\widehat{\mathbf{v}}_{\mathcal{O}}}\leq\norm{\mathbf{v}}$.
\end{enumerate}
\begin{proof}
\begin{enumerate}[label=(\roman*)]
\setlength{\itemsep}{0pt}
\item By construction, $\widehat{\mathbf{v}}_{\mathcal{O}}\in\mathcal{S}$. For $k\in\{1,\ldots,m\}$, $\innerprod{\mathbf{v}-\widehat{\mathbf{v}}_{\mathcal{O}},\mathbf{e}_k}=c_k(\mathbf{v})-\sum\limits_{j=1}^{m}c_j(\mathbf{v})\delta_{jk}=0$.\\
If $\mathbf{x}\in\mathcal{S}$, write $\mathbf{x}=\sum\limits_{k=1}^{m}a_k\mathbf{e}_k$; then $\innerprod{\mathbf{v}-\widehat{\mathbf{v}}_{\mathcal{O}},\mathbf{x}}=\sum\limits_{k=1}^{m}\overline{a_k}\innerprod{\mathbf{v}-\widehat{\mathbf{v}}_{\mathcal{O}},\mathbf{e}_k}=0$ by Proposition \ref{prop:basic-inner-product-identities}(iii)--(iv). Hence $\mathbf{v}-\widehat{\mathbf{v}}_{\mathcal{O}}\in\mathcal{S}^{\perp}$.\\
If $\mathbf{s}\in\mathcal{S}$ also satisfies $\mathbf{v}-\mathbf{s}\in\mathcal{S}^{\perp}$, then $\mathbf{s}-\widehat{\mathbf{v}}_{\mathcal{O}}=(\mathbf{v}-\widehat{\mathbf{v}}_{\mathcal{O}})-(\mathbf{v}-\mathbf{s})\in\mathcal{S}\cap\mathcal{S}^{\perp}$, so $\mathbf{s}=\widehat{\mathbf{v}}_{\mathcal{O}}$ by Proposition \ref{prop:basic-properties-orthogonal-complement}(iv).
\item By (i), $\widehat{\mathbf{v}}_{\mathcal{O}}=P_{\mathcal{S}}\mathbf{v}$. The inequality and equality condition hold by Theorem \ref{thm:orthogonal-projection-minimizes-distance}.
\item If $m=0$, then $\widehat{\mathbf{v}}_{\mathcal{O}}=\mathbf{0}$ and the claim is immediate. Suppose $m\geq1$.\\
By Proposition \ref{prop:orthonormal-linear-combination-norm} and Theorem \ref{thm:bessel-inequality-finite-orthonormal-list}, $\norm{\widehat{\mathbf{v}}_{\mathcal{O}}}^2=\sum\limits_{j=1}^{m}\abs{c_j(\mathbf{v})}^2\leq\norm{\mathbf{v}}^2$, so $\norm{\widehat{\mathbf{v}}_{\mathcal{O}}}\leq\norm{\mathbf{v}}$ by nonnegativity of norms.
\end{enumerate}
\hlt{Structural Commentary}: Coordinate construction and nearest-point optimality prove all three claims.\\
\hlt{Key Insight}: Fourier expansion is orthogonal projection expressed in orthonormal coordinates.
\end{proof}
\end{theorem}

\begin{theorem} \label{thm:finite-orthonormal-basis-characterizations} \hlt{Characterizing Finite Orthonormal Bases}\\
Let $\mathcal{O}=(\mathbf{e}_1,\ldots,\mathbf{e}_m)$ be orthonormal, $\mathcal{S}=\operatorname{span}(\mathcal{O})$. Write $c_j(\mathbf{v})$, $\widehat{\mathbf{v}}\defeq\widehat{\mathbf{v}}_{\mathcal{O}}$ as in Definition \ref{def:fourier-coefficients-and-expansion}. The following are equivalent:
\begin{enumerate}[label=(\roman*)]
\setlength{\itemsep}{0pt}
\item $\mathcal{S}=\mathcal{V}$.
\item $\mathcal{S}^{\perp}=\{\mathbf{0}\}$.
\item $(\forall\,\mathbf{v}\in\mathcal{V})\,\widehat{\mathbf{v}}=\mathbf{v}$.
\item $(\forall\,\mathbf{v},\mathbf{w}\in\mathcal{V})\,\innerprod{\mathbf{v},\mathbf{w}} = \sum\limits_{j=1}^{m}c_j(\mathbf{v})\overline{c_j(\mathbf{w})}$.
\item $(\forall\,\mathbf{v}\in\mathcal{V})\,\norm{\widehat{\mathbf{v}}}=\norm{\mathbf{v}}$.
\end{enumerate}
\begin{proof}
\begin{enumerate}[label=(\roman*)]
\setlength{\itemsep}{0pt}
\item[$(i)$]$\Rightarrow (ii)$ If $\mathcal{S}=\mathcal{V}$, then $\mathcal{S}^{\perp}=\mathcal{V}^{\perp}=\{\mathbf{0}\}$ by Proposition \ref{prop:basic-properties-orthogonal-complement}(iii).
\item[$(ii)$]$\Rightarrow (iii)$ Fix $\mathbf{v}\in\mathcal{V}$. By Theorem \ref{thm:fourier-expansion-best-approximation-finite}(i), $\mathbf{v}-\widehat{\mathbf{v}}\in\mathcal{S}^{\perp}=\{\mathbf{0}\}$, so $\widehat{\mathbf{v}}=\mathbf{v}$. 
\item[$(iii)$]$\Rightarrow (iv)$ Fix $\mathbf{v},\mathbf{w}\in\mathcal{V}$. By (iii), $\mathbf{v}=\sum\limits_{j=1}^{m}c_j(\mathbf{v})\mathbf{e}_j$ and $\mathbf{w}=\sum\limits_{k=1}^{m}c_k(\mathbf{w})\mathbf{e}_k$.\\
Then $\innerprod{\mathbf{v},\mathbf{w}}=\sum\limits_{j=1}^{m}c_j(\mathbf{v})\overline{c_j(\mathbf{w})}$ by Proposition \ref{prop:basic-inner-product-identities}(iii)--(iv) and orthonormality. 
\item[$(iv)$]$\Rightarrow (v)$ Fix $\mathbf{v}\in\mathcal{V}$ and set $\mathbf{w}=\mathbf{v}$ in (iv). Then $\norm{\mathbf{v}}^2=\sum\limits_{j=1}^{m}\abs{c_j(\mathbf{v})}^2=\norm{\widehat{\mathbf{v}}}^2$ by Proposition \ref{prop:orthonormal-linear-combination-norm}, so $\norm{\widehat{\mathbf{v}}}=\norm{\mathbf{v}}$. 
\item[$(v)$]$\Rightarrow (i)$ Fix $\mathbf{v}\in\mathcal{V}$. By Theorem \ref{thm:fourier-expansion-best-approximation-finite}(i), $\widehat{\mathbf{v}}\in\mathcal{S}$ and $\mathbf{v}-\widehat{\mathbf{v}}\in\mathcal{S}^{\perp}$. Hence $\norm{\mathbf{v}}^2=\norm{\widehat{\mathbf{v}}}^2+\norm{\mathbf{v}-\widehat{\mathbf{v}}}^2$ by Theorem \ref{thm:pythagorean-inner-product}.\\
Condition (v) implies $\norm{\mathbf{v}-\widehat{\mathbf{v}}}=0$, so $\mathbf{v}=\widehat{\mathbf{v}}\in\mathcal{S}$ by definiteness. Thus $\mathcal{V}\subseteq\mathcal{S}$, while $\mathcal{S}\subseteq\mathcal{V}$ by construction,  therefore $\mathcal{S}=\mathcal{V}$.
\end{enumerate}
\hlt{Structural Commentary}: A five-step implication cycle proves equivalence.\\
\hlt{Key Insight}: An orthonormal list is a basis exactly when its projection loses no component.
\end{proof}
\end{theorem}

\begin{perspective} \label{per:projection-theorem-finite-dimensional-hypothesis} \hlt{Why the Finite-Dimensional Hypothesis Matters}\\
In $\ell^2(\mathbb{F})$, let $\mathcal{S}=\mathbb{F}^{(\mathbb{N})}=\operatorname{span}\{\mathbf{e}_j:j\geq1\}$. If $\mathbf{x}\in\mathcal{S}^{\perp}$, then $(\forall\,j\in\mathbb{N})\,x_j=\innerprod{\mathbf{x},\mathbf{e}_j}=0$, so $\mathcal{S}^{\perp}=\{\mathbf{0}\}$.\\
The vector $(2^{-j})_{j\geq1}$ belongs to $\ell^2(\mathbb{F})$ but not to $\mathcal{S}$, hence $\mathcal{S}\subsetneq\ell^2(\mathbb{F})$ and $\mathcal{S}\oplus\mathcal{S}^{\perp}\neq\ell^2(\mathbb{F})$.
\end{perspective}
